\documentclass[10pt,a4paper]{amsart}
\usepackage[margin=19mm]{geometry}
\usepackage{amsmath,amssymb,mathtools}
\usepackage[scr=rsfs,cal=euler]{mathalfa}
\usepackage{xfrac}
\usepackage[unicode,hidelinks,bookmarks=false]{hyperref}
\newcommand{\ie}{i.e.\ }
\newcommand{\cf}{cf.\ }
\newcommand{\resp}{respectively\ }
\newcommand{\Subsection}{\subsection}
\providecommand{\nobreakdash}{\nobreak\hbox{-}}
\setlength{\emergencystretch}{2em}
\multlinegap0pt
\usepackage{braket}
\usepackage{amssymb}
\usepackage{bm}
\usepackage{tikz}
\newtheorem{lemma}{Lemma}
\newcommand{\CC}{\mathsf{C}\mathsf{C}}
\newcommand{\rI}{\mathsf{I}}
\newcommand{\rR}{\mathsf{R}}
\newcommand{\rU}{\mathsf{U}}
\title{A Rigorous and Self-Contained Proof of the Grover–Rudolph State Preparation Algorithm}
\author{Antonio Falc\'o}
\date{Source: arXiv:2601.17930v2 (CC BY 4.0)}
\begin{document}
\maketitle

\noindent\emph{Source and licence.} A Rigorous and Self-Contained Proof of the Grover–Rudolph State Preparation Algorithm. Version arXiv:2601.17930v2 (CC BY 4.0), distributed by arXiv under the Creative Commons Attribution 4.0 International licence, http://creativecommons.org/licenses/by/4.0/, as recorded in the arXiv licence metadata for this exact version and shown on the abstract page https://arxiv.org/abs/2601.17930v2. Full licence notice: \texttt{licenses/arxiv-2601.17930-CC-BY-4.0.txt}.\par\smallskip

\noindent\textit{Editorial scope.} Counted unit: the article's lemma lem:inductive\_state (source0.tex lines 778--797). The article's complete original proof of it is delivered with it (source0.tex lines 799--930, one \texttt{proof} environment of 132 lines). No mathematics is written, repaired or re-derived: the statement and the proof are the article's own lines, in the article's own notation, and no retained line is substituted or re-flowed.  Retained uncounted as the source-local context the proof needs: lines 519--529.  Nothing was omitted from any retained span, so the package carries no omission record.  No retained line contains a citation, so the wrapper carries no bibliography and no bibliography entry.  No retained line contains a figure environment, an image-inclusion command or drawing code, so no image is delivered and the package declares no asset dependency.  Editorial formatting only, in addition: an AMS \texttt{amsart} preamble that restates the article's own declarations and macros used by the retained lines, namely the environments \texttt{lemma} on separate counters, together with the standard packages those lines need.  The retained environments print in this document as Lemma 1 in order of appearance, whereas the article prints the same items with the numbering of its own full document.  The complete native source of the article as distributed in the arXiv e-print for this exact version is bundled unchanged as \texttt{sources/193351/source0.tex}.
\begin{lemma}[Suffix extraction via scaling]\label{lem:suffix-extraction}
Let $k\in \mathbb{Z}_{2^n}$ with $b_n(k)=z_1z_2\cdots z_n$.
For $1\le s\le n-1$ we have
\[
\Big\lfloor 2^{-s}k\Big\rfloor
=
\sum_{j=s+1}^n z_j\,2^{j-s-1},
\qquad\text{and hence}\qquad
b_{n-s}\!\Big(\lfloor 2^{-s}k\rfloor\Big)=z_{s+1}\cdots z_n.
\]
\end{lemma}

\begin{lemma}[Inductive amplitude form]\label{lem:inductive_state}
Let $\ket{\psi_j}:=\rU_j\rU_{j-1}\cdots \rU_1\ket{0}^{\otimes n}$.
Then, for every $1\le j\le n$,
\[
\ket{\psi_j}
=
\ket{0}^{\otimes (n-j)}\otimes
\sum_{\mathbf{w}=z_{n-j+1}\cdots z_n\in\mathbb{Z}_2^{j}}
\Bigl(\mathsf{T}_{z_{n-j+1}}(\theta_{z_{n-j+2}\cdots z_n})\cdots
\mathsf{T}_{z_{n-1}}(\theta_{z_n})\mathsf{T}_{z_n}(\theta)\Bigr)\,
\ket{\mathbf{w}}.
\]
In particular,
\[
\ket{\psi_n}
=
\sum_{z_1\cdots z_n\in\mathbb{Z}_2^n}
\Bigl(\mathsf{T}_{z_1}(\theta_{z_2\cdots z_n})\cdots \mathsf{T}_{z_n}(\theta)\Bigr)\ket{z_1\cdots z_n}.
\]
\end{lemma}

\begin{proof}
We proceed by induction on $j$.

\smallskip
\noindent\textbf{Base case ($j=1$).}
By definition, $\rU_1=\rI_{2^{n-1}}\otimes \rR(\theta)$ acts only on the last qubit. Hence
\[
\ket{\psi_1}
=\rU_1\ket{0}^{\otimes n}
=\ket{0}^{\otimes(n-1)}\otimes \rR(\theta)\ket{0}.
\]
Since $\rR(\theta)\ket{0}=\cos\theta\,\ket{0}+\sin\theta\,\ket{1}$, we obtain
\[
\ket{\psi_1}
=\ket{0}^{\otimes(n-1)}\otimes\bigl(\cos\theta\,\ket{0}+\sin\theta\,\ket{1}\bigr)
=\ket{0}^{\otimes(n-1)}\otimes\sum_{z_n\in\mathbb{Z}_2}\mathsf{T}_{z_n}(\theta)\ket{z_n},
\]
which is the claimed formula for $j=1$.

\smallskip
\noindent\textbf{Inductive step.}
Assume the statement holds for some $j-1$ with $2\le j\le n$, i.e.
\begin{equation}\label{eq:psi_jminus1}
\ket{\psi_{j-1}}
=
\ket{0}^{\otimes (n-(j-1))}\otimes
\sum_{\mathbf{w}=z_{n-j+2}\cdots z_n\in\mathbb{Z}_2^{j-1}}
A_{\mathbf{w}}\,\ket{\mathbf{w}},
\end{equation}
where, for each $\mathbf{w}=z_{n-j+2}\cdots z_n$, the amplitude $A_{\mathbf{w}}$ is given by the product
\begin{equation}\label{eq:Aw_def}
A_{\mathbf{w}}
=
\mathsf{T}_{z_{n-j+2}}(\theta_{z_{n-j+3}\cdots z_n})\cdots
\mathsf{T}_{z_{n-1}}(\theta_{z_n})\,
\mathsf{T}_{z_n}(\theta).
\end{equation}

At this stage, the basis vectors $\ket{\mathbf{w}}$ label the computational states of the last $(j-1)$ qubits.
Their identification with suffixes of the global binary index follows from the fact that truncating
the first $(j-1)$ bits of an integer corresponds to extracting its suffix.
More precisely, if $k\in\mathbb{Z}_{2^n}$ has binary expansion $b_n(k)=z_1\cdots z_n$, then
\[
b_{n-(j-1)}\!\bigl(\lfloor 2^{-(j-1)}k\rfloor\bigr)=z_{n-j+2}\cdots z_n,
\]
as shown in Lemma~\ref{lem:suffix-extraction}.



It will be convenient to isolate the $(n-j+1)$-th qubit (the next qubit to be rotated at stage $j$).
Write the tensor product in \eqref{eq:psi_jminus1} as
\[
\ket{0}^{\otimes(n-(j-1))}
=\ket{0}^{\otimes(n-j)}\otimes \ket{0},
\]
so that
\begin{equation}\label{eq:psi_jminus1_rewrite}
\ket{\psi_{j-1}}
=
\ket{0}^{\otimes(n-j)}\otimes
\sum_{\mathbf{w}\in\mathbb{Z}_2^{j-1}} A_{\mathbf{w}}\,\bigl(\ket{0}\otimes\ket{\mathbf{w}}\bigr),
\end{equation}
where the factor $\ket{0}$ corresponds to qubit $(n-j+1)$ and $\ket{\mathbf{w}}$ to the last $(j-1)$ qubits.

\smallskip
\noindent\emph{Action of $\rU_j$ on basis branches.}
By definition,
\[
\rU_j=\prod_{\mathbf{w}\in\mathbb{Z}_2^{j-1}}\Bigl(\rI_{2^{n-j}}\otimes \CC_{\mathbf{w}}^{(1)}\bigl(\rR(\theta_{\mathbf{w}})\bigr)\Bigr).
\]
Inside the active $j$-qubit register (qubits $(n-j+1),\dots,n$), the gate
$\CC_{\mathbf{w}}^{(1)}(\rR(\theta_{\mathbf{w}}))\in\mathrm{U}(2^j)$ acts as follows:
it applies $\rR(\theta_{\mathbf{w}})$ to the \emph{first} qubit of that register (i.e.\ qubit $(n-j+1)$)
if and only if the remaining $(j-1)$ qubits are in the computational basis state $\ket{\mathbf{w}}$, and
acts as the identity on all other computational basis states.
Consequently, for each $\mathbf{w}\in\mathbb{Z}_2^{j-1}$,
\begin{equation}\label{eq:control_action}
\CC_{\mathbf{w}}^{(1)}\bigl(\rR(\theta_{\mathbf{w}})\bigr)\,\bigl(\ket{0}\otimes\ket{\mathbf{w}}\bigr)
=\bigl(\rR(\theta_{\mathbf{w}})\ket{0}\bigr)\otimes\ket{\mathbf{w}},
\end{equation}
and for $\mathbf{w}'\neq \mathbf{w}$,
\[
\CC_{\mathbf{w}}^{(1)}\bigl(\rR(\theta_{\mathbf{w}})\bigr)\,\bigl(\ket{0}\otimes\ket{\mathbf{w}'}\bigr)
=\ket{0}\otimes\ket{\mathbf{w}'}.
\]
Since the projectors onto distinct control branches are orthogonal, these controlled rotations act on
disjoint subspaces and therefore commute; hence the product over $\mathbf{w}$ applies the correct rotation in each
branch independently.

\smallskip
\noindent\emph{Splitting of amplitudes.}
Applying $\rU_j$ to \eqref{eq:psi_jminus1_rewrite} and using \eqref{eq:control_action} gives
\begin{align*}
\ket{\psi_j}
&=\rU_j\ket{\psi_{j-1}}\\
&=
\ket{0}^{\otimes(n-j)}\otimes
\sum_{\mathbf{w}\in\mathbb{Z}_2^{j-1}} A_{\mathbf{w}}\,
\Bigl(\rR(\theta_{\mathbf{w}})\ket{0}\Bigr)\otimes\ket{\mathbf{w}}.
\end{align*}
Finally, $\rR(\theta_{\mathbf{w}})\ket{0}=\cos\theta_{\mathbf{w}}\,\ket{0}+\sin\theta_{\mathbf{w}}\,\ket{1}
=\sum_{z\in\mathbb{Z}_2}\mathsf{T}_{z}(\theta_{\mathbf{w}})\ket{z}$, so
\begin{align*}
\ket{\psi_j}
&=
\ket{0}^{\otimes(n-j)}\otimes
\sum_{\mathbf{w}\in\mathbb{Z}_2^{j-1}} A_{\mathbf{w}}\,
\sum_{z_{n-j+1}\in\mathbb{Z}_2}\mathsf{T}_{z_{n-j+1}}(\theta_{\mathbf{w}})\,
\ket{z_{n-j+1}}\otimes\ket{\mathbf{w}}\\
&=
\ket{0}^{\otimes(n-j)}\otimes
\sum_{z_{n-j+1}\in\mathbb{Z}_2}\;
\sum_{\mathbf{w}\in\mathbb{Z}_2^{j-1}}
\Bigl(A_{\mathbf{w}}\,\mathsf{T}_{z_{n-j+1}}(\theta_{\mathbf{w}})\Bigr)\,
\ket{z_{n-j+1}\mathbf{w}}.
\end{align*}
Renaming the concatenated word $z_{n-j+1}\mathbf{w}$ as
\[
\mathbf{w}'=z_{n-j+1}z_{n-j+2}\cdots z_n\in\mathbb{Z}_2^j,
\]
and substituting $A_{\mathbf{w}}$ from \eqref{eq:Aw_def}, we obtain exactly the claimed product form
\[
\ket{\psi_j}
=
\ket{0}^{\otimes (n-j)}\otimes
\sum_{\mathbf{w}'=z_{n-j+1}\cdots z_n\in\mathbb{Z}_2^{j}}
\Bigl(\mathsf{T}_{z_{n-j+1}}(\theta_{z_{n-j+2}\cdots z_n})\cdots
\mathsf{T}_{z_{n-1}}(\theta_{z_n})\mathsf{T}_{z_n}(\theta)\Bigr)\,
\ket{\mathbf{w}'}.
\]
This completes the induction. The final statement for $\ket{\psi_n}$ is the special case $j=n$.
\end{proof}

\end{document}
